\begin{theorem}[Algorithm 5.5 base stage]
Let $\mathcal E$ be a positive-rung ETT with rung number $n$, terminal
ladder tree $T_n$, final series output $T_{n+1}$, generating coloring
$\varphi_n$, and maximum property.  Assume also that $\varphi_n$ is a
\noderef{ModColoring} of $\mathcal E$.  Fix a
\noderef{HierarchyBase} $B$.

In the PE case assume
\[
 |T_n|<|B.\mathit{core}|,
 \qquad
 \noderef{TashkinovClosure}
   (\varphi_n,B.\mathit{core},T_{n+1}),                         \tag{1}
\]
and the following three facts from (5.3)--(5.4): some color missing on
$B.\mathit{core}$ is not closed there; every color in
\[
 C=\overline\varphi_n(T_n)\cap
   \overline\varphi_n(B.\mathit{closure})
\]
is closed on $B.\mathit{core}$; and $|C|\geq4$.  Finally assume the
large-branch estimate
\[
                    |\overline\varphi_n(T_n)|\geq2n+11.       \tag{2}
\]

Then there are an exact Algorithm 5.5 seed $J$, a terminal ordered tree
$U$, and a \noderef{CorrespondingExtendedTashkinovTree} $\mathcal E_U$
for $(\varphi_n,T_n)$ on $U$ with the following properties.
\begin{enumerate}
\item In PE, $J=B.\mathit{core}=T_n\vee R_n$.  In SE or RE, there is an
  individual edge identity $f_n$ for which $F_n=\{f_n\}$ and
  \noderef{TreeSequenceAdjoinEdge}$(T_n,f_n,J)$ holds.
\item There is a $\varphi_n$-\noderef{TashkinovAugmentationSequence}
  from $J$ to $U$, and $J$ is an exact ordered initial segment of $U$.
  In PE, $|B.\mathit{core}|<|U|$; in the non-PE cases, $|T_n|<|U|$.
\item The corresponding ETT $\mathcal E_U$ has maximum property and a
  hierarchy base $B_U$ whose core, closure, and numerical support position
  equal those of $B$.  It carries a zero-divider hierarchy $K_U$ whose
  exact stages are $T_n$ and $U$.
\item There is a good-hierarchy witness $\mathcal G_U$ on $K_U$.  At its
  sole indexed stage, every stored color $\eta$ has a two-color set
  $\Gamma^0_\eta$ contained in
  \[
   \overline\varphi_n(T_n)-
     \noderef{HierarchySegmentColors}(K_U,0,\eta),             \tag{3}
  \]
  and the pairs belonging to distinct stored colors are disjoint.  In PE,
  $B.\mathit{core}$ is not $(\Gamma^0)^-$-closed and
  $|C-\Gamma^0|\geq4$, where
  $\Gamma^0=\bigcup_{\eta\in D^{K_U}_{n,0}}\Gamma^0_\eta$.
\item With the terminal set kept separate from the sole indexed set,
  $U$ is
  \[
   \left(\bigcup_{\eta\in
      \noderef{HierarchyTerminalStoredColors}(K_U)}
      \Gamma^0_\eta\right)^-
  \text{-closed under }\varphi_n.                              \tag{4}
  \]
\end{enumerate}
Every stage equality and initial-segment relation above concerns the full
ordered edge-valued tree-sequence.  In particular, item 1 retains the exact
identity of $f_n$ in the SE and RE cases.
\end{theorem}
\begin{proof}
Apply \noderef{AlgorithmFiveFiveSeed} with the maximum property, the base,
and the two conditional assumptions in (1).  It returns an exact seed
$J$, an identity corresponding presentation $\mathcal E_J$, maximum
property for $\mathcal E_J$, a copied base $B_J$, and a zero-divider
hierarchy $K_J$.  Its statement gives
\[
 B_J.\mathit{core}=B.\mathit{core},\quad
 B_J.\mathit{closure}=B.\mathit{closure},quad
 K_J(0)=T_n,quad K_J(1)=J,                                   \tag{5}
\]
as exact data.  It also gives exactly the two alternatives in item 1.  In
the non-PE alternative it retains an individual $f_n$ with
$F_n=\{f_n\}$ and $T_n+f_n=J$; no endpoint equality is used in place of
this edge identity.

Put $D=D^{K_J}_{n,0}=\noderef{HierarchyStoredColorsAt}(K_J,0)$ and
$M=\overline\varphi_n(T_n)$.  We next choose the initial Gamma pairs.  If
the terminal kind is PE, transport the three PE assumptions across the
first two equalities in (5), and transport (2) across the ladder equality
in the corresponding certificate.  The hypotheses of
\noderef{GoodHierarchyInitialPairsPE} are now exact.  It returns a family
$\Gamma_\eta$ such that
\[
 |\Gamma_\eta|=2,quad \Gamma_\eta\subseteq M,quad
 \Gamma_\eta\cap\Gamma_\theta=\varnothing
       \quad(\eta\ne\theta\text{ in }D),                       \tag{6}
\]
the PE core is not $(\bigcup_{\eta\in D}\Gamma_\eta)^-$-closed,
and at least four colors of $C$ remain outside this union.

If the terminal kind is not PE, the positive terminal certificate in
\noderef{TashkinovSeries} and \noderef{TashkinovSeriesStep} must instead
be SE or RE.  Apply
\noderef{GoodHierarchyInitialPairsNonPE}, again using the corresponding
ladder equality to rewrite (2).  It gives exactly the three assertions in
(6).  Thus a family satisfying (6) is fixed in every case.

We record one fact needed for the non-PE avoidance check.  In an SE
terminal certificate, the exact edge $f_n$ is the edge of its
\noderef{MaximumDefectiveChoice}; its selected coloring gives
$\varphi_n(f_n)=\delta_n$, and it is a boundary edge of $T_n$.  The
pre-step coloring is \noderef{ColoredVertexSetClosed} on $T_n$, while the
selected coloring is stable on $T_n$.  Both are proper and leave exactly the distinguished edge
uncolored.  Symmetry of stability from \noderef{sc2} transports a
hypothetical selected-color boundary edge whose color is selected-missing
back to a pre-step boundary edge with the same color and the same missing
witness.  This contradicts pre-step closedness.  Hence $T_n$ is closed
under $\varphi_n$, and the displayed boundary edge gives
\[
                         \delta_n\notin M.                    \tag{7}
\]
In an RE terminal certificate, the current coloring equals the preceding
closed coloring.  The statement of
\noderef{RevisitingExtensionCertificate} says that its exact edge $f_n$
is a boundary edge and lies on the recorded \noderef{TwoColorCycle}, so it has one of the two cycle colors under
$\varphi_n$.  Closedness again shows directly that this exact edge color
is not in $M$.  Consequently, in either non-PE case,
\[
                  \varphi_n(f_n)=c_n\quad\text{for some }c_n\notin M.
                                                                    \tag{8}
\]

Define a total Gamma family on $K_J$ by the selected sets
$\Gamma_\eta$.  We verify the five fields of
\noderef{GoodHierarchy}.  Pair size and pairwise disjointness are (6).
At index zero \noderef{HierarchyStageMissingColors} is $M$.  In PE, (5)
and \noderef{HierarchyStarredStage} say that the starred stage and
terminal seed are both $B.\mathit{core}$, so the
half-open edge interval in
\noderef{HierarchySegmentColors} is empty.  Thus
$\Gamma_\eta\subseteq M$ proves avoidance.  In SE or RE, the starred
stage is $T_n$, the terminal seed is the exact one-edge tree $T_n+f_n$,
and the only possible segment color is the color $c_n$ of that exact new
edge.  Equations (6) and (8) show that no Gamma pair contains $c_n$, so
avoidance holds here as well.

There is no positive divider index, so `uniqueRefresh' and the internal
`minusClosed' field have empty domains.  In PE the two parts of
`parallelBase' are precisely the non-minus-closed and four-color
conclusions returned by \noderef{GoodHierarchyInitialPairsPE}, after the
equalities in (5).  Outside PE its antecedent is false.  We have therefore
constructed a complete good witness $\mathcal G_J$ on $K_J$.

Now apply \noderef{GoodHierarchyRestrictedAugmentation} to
$\mathcal E$, its mod-coloring and maximum property, and the corresponding
seed data $(J,\mathcal E_J,B_J,K_J,\mathcal G_J)$.  It supplies a finite
chain
\[
                         J=U_0\prec U_1\prec\cdots\prec U_m=U, \tag{9}
\]
where every nonempty step is an exact
$\varphi_n$-Tashkinov augmentation, together with its composite
augmentation sequence and exact initial-segment relations.  It also
supplies a corresponding ETT $\mathcal E_U$ with maximum property, a base
$B_U$, a hierarchy $K_U$, and a good witness $\mathcal G_U$.

The reconstruction clauses of that lemma identify the core, closure, and
numerical support position of $B_U$ with those of $B_J$, and hence with
those of $B$.  They identify the divider count of $K_U$ with the zero
divider count of $K_J$, identify its sole nonterminal stage with $T_n$,
and identify every indexed stored set and Gamma pair with the corresponding
one in $K_J$.  The terminal-stage field of
\noderef{TashkinovHierarchy} identifies its other stage with the exact
tree $U$.  Thus item 3, (3), pair cardinality, pairwise disjointness, and
the two PE clauses all transport literally from (5)--(6) and
$\mathcal G_J$ to $K_U$ and $\mathcal G_U$.

We prove the strict assertions in item 2.  In SE or RE,
\noderef{TreeSequenceAdjoinEdge} gives $|J|=|T_n|+1$, and (9) gives
$J\preccurlyeq U$, so $|T_n|<|U|$.

Suppose the terminal kind is PE.  Let
$R=\bigcup_{\eta\in D}\Gamma_\eta$.  The PE selection says that the seed
$J=B.\mathit{core}$ is not $R^-$-closed.  Expanding
\noderef{TreeColorSetMinusClosed}, choose a color
$c\in\overline\varphi_n(J)-R$ and an exact $c$-colored boundary edge $f$
of $J$.  By \noderef{BoundaryColorEdges} and \noderef{BoundaryEdges}, it
has an inside endpoint and an outside endpoint with the asserted unordered
endpoint pair.  These endpoints construct the ordered one-edge tree
$W=J+f$: in \noderef{TreeSequence}, copy every old vertex and edge entry,
append the outside endpoint and the individual identity $f$, and use the
boundary endpoint equation for the final attachment.  Freshness of the
outside endpoint proves vertex injectivity and also prevents $f$ from
being an old displayed edge, proving edge injectivity.  The last edge has color $c$
missing on $J$, so this is a
\noderef{TashkinovAugmentation} from $J$ to $W$.

This step satisfies every fixed Gamma-avoidance condition used by the
restricted loop.  All old eligible segment colors were already disjoint
from $\Gamma_\eta$ by the avoidance field of $\mathcal G_J$; the only new
possible segment color is $c$, and $c\notin R$ excludes it from every
$\Gamma_\eta$.  Hence $W$ is an admissible successor of $J$.

If the terminal tree in (9) were $U=J$, the maximality clause returned by
\noderef{GoodHierarchyRestrictedAugmentation} would forbid exactly this
successor $W$, a contradiction.  Thus $U\ne J$.  The length formula in
that lemma says $|U|=|J|+m$, so $m>0$ and $|J|<|U|$.  Since $J$ is the PE
core, this is the required strict PE growth.

Finally, the terminal conclusion of
\noderef{GoodHierarchyRestrictedAugmentation} says directly that
\[
 \noderef{TreeColorSetMinusClosed}\left(
   G,\varphi_n,U,
   \bigcup_{\eta\in D^{K_U}_{n,1}}
     \mathcal G_U.\Gamma^0_\eta\right)                         \tag{10}
\]
holds, where its displayed terminal set is exactly
\noderef{HierarchyTerminalStoredColors}$(K_U)$.  Equation (10) is item 5
and (4).  All requested seed, growth, correspondence, maximum, hierarchy,
avoidance, disjointness, PE, and terminal closedness assertions have now
been established.
\end{proof}
